
% --- SLIDE 5: CI/CD ---
\begin{frame}{CI/CD for ML: What Runs When}

    \centering
    \renewcommand{\arraystretch}{1.3}
    {\small
    \begin{tabular}{p{3.0cm}|p{4.2cm}|p{4.2cm}}
        \toprule
        \textbf{Trigger} & \textbf{What runs} & \textbf{Gate} \\
        \midrule
        Every commit / PR & Linting, unit tests, a training run on a \emph{tiny} sample &
            Must be green to merge; must finish in minutes \\
        Merge to main     & Build the container, publish it by digest &
            Image builds and the smoke test passes \\
        New data or a schedule & Full training pipeline, with data validation first &
            Candidate beats the incumbent on the frozen evaluation set \\
        Manual promotion  & Register the model version, then staged rollout &
            A human approves --- or an automated policy does, with a rollback plan \\
        \bottomrule
    \end{tabular}}

    \vspace{0.9em}
    \begin{itemize}
        \item \textbf{Keep the per-commit job fast.} A CI job that trains the real model for six hours is a
              CI job people learn to skip.
        \item \textbf{Continuous training} (retraining on fresh data) and \textbf{continuous delivery}
              (shipping new pipeline code) are different loops running at different cadences. Do not
              entangle them.
    \end{itemize}

\end{frame}
